\section*{Chapitre \ref{ch:m3:centralised-exchanges} --- Plateformes centralisées}

\begin{solution}{exo:m3:centralised-exchanges:1}
Le courtier (le compte), la bourse (l'appariement), la contrepartie centrale (garantie et marge) et
le dépositaire ou le système de règlement-livraison (qui détient les actifs).
\end{solution}

\begin{solution}{exo:m3:centralised-exchanges:2}
2\,400 dans les deux cas~: 1\,200 juste avant une frontière de fenêtre et 1\,200 juste après
(\cref{prop:m3:centralised-exchanges:burst})~; aucun intervalle d'une minute ne peut recouper plus de deux fenêtres.
\end{solution}

\begin{solution}{exo:m3:centralised-exchanges:3}
La clé d'une stratégie réside sur des serveurs et dans une configuration que de nombreuses personnes et de nombreux programmes manipulent~; en cas de fuite,
une clé réservée au trading permet à un attaquant de passer des ordres préjudiciables, une clé de retrait lui permet de s'emparer des actifs. Les retraits passent
par une procédure distincte, avec approbations et adresses sur liste blanche.
\end{solution}

\begin{solution}{exo:m3:centralised-exchanges:4}
Niveau standard~: $25\text{ millions} \times 0.10\% \times 2 = \$50\,000$ par jour. VIP~9~: $25\text{ millions} \times
(0.011\% + 0.023\%) = \$8\,500$ par jour.
\end{solution}

\begin{solution}{exo:m3:centralised-exchanges:5}
L'arbre compte 21 niveaux au-dessus des feuilles ($2^{21} = 2\,097\,152$), donc 21 hachages après celui de sa propre feuille.
$1 - 0.99^{100} = 63.4\%$.
\end{solution}

\begin{solution}{exo:m3:centralised-exchanges:6}
$(270 + 310)/6\,991 = 8.3\%$.
\end{solution}

\begin{solution}{exo:m3:centralised-exchanges:7}
15 clients le détectent (1\,008 passent). Une fausse feuille dissimulant 2~\% ne rend négatifs que ses ancêtres jusqu'au
sous-arbre de 16 feuilles, dont les soldes réels dépassent ce montant~; seuls les clients dont le chemin a pour frère l'un de ces nœuds
négatifs ($1 + 2 + 4 + 8 = 15$) voient une somme négative. Dissimuler moins rend la détection plus difficile.
\end{solution}

\begin{solution}{exo:m3:centralised-exchanges:8}
Une \omterm{def:m3:centralised-exchanges:por}{preuve de réserves} est un instantané~: elle ne montre ni les passifs envers les prêteurs, ni les actifs empruntés pour l'instantané,
ni ceux prêtés après coup, ni les clés que d'autres détiennent aussi, ni tout ce qui s'est passé depuis. Et même une plateforme solvable est une
contrepartie non garantie~: 40 millions de dollars y constituent une exposition de crédit à limiter et à transférer, et non une position «~sûre~».
\end{solution}

\begin{solution}{pb:m3:centralised-exchanges:1}
\textbf{1.} 7\,014 millions de dollars (total du dossier~; 7\,013 à partir des lignes quotidiennes arrondies). \textbf{2.} 7
novembre~: 3\,306 millions de dollars. \textbf{3.} Il y a eu 3\,241 millions de dollars de dépôts nets, surtout entre le 6 et le 8 novembre~;
le dossier précise que la forte hausse des dépôts provient en grande partie d'Alameda. Les dépôts de l'entité affiliée
ont gonflé les soldes inscrits au livre de la plateforme pendant que les clients retiraient leurs actifs réels. \textbf{4.} La plateforme
n'avait plus d'actifs pour payer et avait en pratique cessé de traiter les retraits. \textbf{5.} Les
clients qui ont retiré les premiers, tant qu'il restait des actifs~: les retraits se font dans l'ordre d'arrivée. \textbf{6.} $694/10\,544
= 6.6\%$. \textbf{7.} $1\,078/10\,544 = 10.2\%$. \textbf{8.} $1/1\,591 = 0.06\%$ localisés, 0.4~\% avec
créances. \textbf{9.} $2\,540/11\,233 = 22.6\%$~; un déficit de 8\,693 millions de dollars. \textbf{10.} Elle compte
des \omterm{def:m3:blockchains-for-traders:key}{jetons} peu liquides liés au groupe valorisés aux cours affichés (MAPS~: 1\,004 millions localisés contre 96 dus) qui
ne pouvaient être vendus nulle part à proximité de ces cours. \textbf{11.} Seulement en partie~: elle aurait montré à chaque titulaire d'un compte à vue
que son solde était pris en compte, et non que les actifs correspondaient au total. \textbf{12.} Oui, si la comparaison
était honnête et complète~: un million de dollars en bitcoins face à 1\,591 millions dus ne peut être dissimulé, sauf si
des actifs avaient été empruntés pour l'instantané. \textbf{13.} Non~: le privilège tenait aux réglages du moteur
d'appariement, et non aux soldes. \textbf{14.} Le collatéral confié à un dépositaire indépendant serait resté hors
de la masse~; la perte aurait porté sur les gains non réglés à la plateforme. \textbf{15.} Une limite de crédit dimensionnée
aux fonds propres de la société et à la qualité de la plateforme, sur la valeur de marché des soldes majorée des retraits en attente
et des plus-values latentes, l'excédent étant transféré chaque jour. \textbf{16.} Le Bankruptcy Code fixe les créances en dollars
à la date de la requête~; les clients, créanciers d'actifs qui ont ensuite progressé, ont perdu cette hausse, que la masse a conservée.
\textbf{17.} Les actifs ségrégués des clients restent hors de la masse d'un courtier défaillant~; un solde sur une plateforme, en l'absence
de telles clauses et de tels textes, est une créance non garantie. \textbf{18.} Les doutes publics sur le bilan de l'entité affiliée,
la concentration dans le \omterm{def:m3:blockchains-for-traders:key}{jeton} propre au groupe, la baisse de son cours, des retraits plus lents~: autant de signes incitant à réduire les soldes
avant que les faits ne soient établis avec certitude. \textbf{19.} \emph{Résultat nommé~:} 6.6~\% des créances couvertes par
des \omterm{def:m3:blockchains-for-traders:key}{jetons} liquides localisés, 10.2~\% avec créances. \textbf{20.} Un prêt non garanti à la plateforme, à moins que ses
conditions et la loi n'en fassent un prêt mieux garanti.
\end{solution}

\begin{solution}{iq:m3:centralised-exchanges:1}
Un ordre à cours limité qui est rejeté (ou dont le cours est ajusté) s'il devait s'exécuter aussitôt, de sorte qu'il ne fasse que reposer au carnet~: il
assure le tarif maker et empêche qu'une cotation calculée à partir d'un prix périmé ne croise la fourchette. Les teneurs
de marché et les stratégies d'exécution passive y recourent.
\iqlookfor{tarif maker, aucune prise involontaire et le rejet comme issue attendue.}
\end{solution}

\begin{solution}{iq:m3:centralised-exchanges:2}
Il faut traiter les soldes comme un crédit non garanti~: limites par plateforme dimensionnées selon les fonds propres et la qualité de la plateforme, soldes limités
à ce qu'exige le trading et transférés vers une conservation propre ou tierce, \omterm{def:m3:centralised-exchanges:oes}{règlement hors plateforme} là où il est proposé, \omterm{def:m3:centralised-exchanges:subaccount}{sous-comptes}
pour séparer les stratégies, surveillance des signes de tension (délais de retrait, \omterm{def:m3:blockchains-for-traders:key}{jeton} de la plateforme, actualités) assortie d'un plan
pour retirer rapidement les fonds.
\iqlookfor{limites de crédit, rapatriement des soldes et conservation hors de la plateforme.}
\end{solution}

\begin{solution}{iq:m3:centralised-exchanges:3}
Il faut y voir un bug~: s'arrêter et respecter \texttt{Retry-After}, car de nouvelles requêtes entraîneraient un bannissement IP~; puis
chercher pourquoi la comptabilité des poids et des ordres tenue par le client a divergé de celle de la plateforme (en-têtes d'usage, poids
modifiés, plusieurs processus partageant une même IP), corriger le limiteur et regrouper les mises à jour pour que les rafales ne
dépassent pas les limites.
\iqlookfor{attente progressive, comptabilité côté client et bannissement.}
\end{solution}

\begin{solution}{iq:m3:centralised-exchanges:4}
Elle reçoit le nonce de sa feuille, les hachés frères et les sommes le long de son chemin, recalcule la racine et la compare
à celle qui est publiée, en vérifiant que chaque somme est positive ou nulle. La plateforme pourrait omettre les comptes peu susceptibles
d'être contrôlés, ou ajouter des soldes négatifs pour réduire le total si les vérificateurs ne contrôlent pas les sommes.
\iqlookfor{le chemin, les sommes et les deux fraudes possibles.}
\end{solution}

\begin{solution}{iq:m3:centralised-exchanges:5}
Les passifs autres que les soldes des clients, les actifs empruntés pour l'instantané ou prêtés ensuite, le contrôle exclusif
des clés, les actifs et passifs hors chaîne, et tout ce qui survient après l'instantané.
\iqlookfor{un instantané d'un seul côté du bilan.}
\end{solution}

\begin{solution}{iq:m3:centralised-exchanges:6}
400 cotations (deux côtés) ne peuvent pas toutes être mises à jour à 10 ordres par seconde~; conserver pour chaque instrument et chaque côté la cotation souhaitée et
celle en place, classer les modifications en attente par valeur (distance au prix juste, instruments en
mouvement), et laisser le limiteur libérer les modifications les plus utiles dans la limite autorisée, en regroupant pour que seule la
dernière cotation souhaitée soit envoyée~; élargir ou retirer les cotations qu'il est impossible de tenir à jour, et utiliser l'amendement ou l'annulation-remplacement
selon leur décompte par la plateforme.
\iqlookfor{regroupement, hiérarchisation et limiteur.}
\end{solution}
